
The quality of MRP inference depends on how well the regression is able to model opinion within each cell.
Researchers can sometimes assess error by comparing model-based estimates to known quantities --- typically in the context of election surveys.
We do not know how particular individuals voted, but we do know election outcomes at various levels of geographic aggregation. 
Even if predicting opinion in these geographies is not of primary interest, these discrepancies can be used to  improve inference for other subgroups. 

A popular method of adjusting MRP estimates to account for such error involves adding a geography-specific intercept shift to the modeled probabilities in each poststratification cell.\footnote{In Appendix~\ref{apx:otherwting}, we review how to account for this information in a classic survey weighting context and discuss limitations to that approach.}  
The values of these geographic intercepts are chosen such that the model-implied vote shares exactly match the known population vote shares in each geographic unit \citep[769]{GG:2013}.
Because the intercept adjustment is applied on the logit scale, this method has been called the ``logit shift.''

To formalize this method, suppose we observe the population-level outcome $\theta_g^j$ for some geography $g$. 
Given model estimates $\hat{\theta}_g^j$, we calibrate results by adding an intercept shift on the logit scale to the predicted probabilities.
The logit shift parameter for geography $g$ on outcome $j$ is the value of $\delta_g^j$ that solves:
\begin{align}
	\underbrace{\theta_g^j}_{\substack{\text{ground truth} \\ \text{for geography $g$}}} &= \,\, \underbrace{\frac{1}{N_g} \sum_{c \in \mathcal{C}_g} N_c \,\, \overbrace{\text{logit}^{-1}\left(\delta_{g[c]}^j + \text{logit}(\hat{\theta}_c^j) \right)}^{\text{updated estimate for cell $c$}}}_{\text{updated estimate for geography $g$}}, \label{eq:logitshift}
\end{align}
where $N_g \equiv \sum_{c \in \mathcal{C}_g} N_c$ is the population of geography $g$ and the notation $g[c]$ indicates the geography for cell $c$.
This expression equates the adjusted model-implied outcome to the ground-truth population outcome $\theta_g^j$. 

The new cell-level probabilities are generated by applying the geographic offset to the existing probabilities: 
\begin{align}
	\tilde{\theta}_c^j = \text{logit}^{-1}\left( \delta_{g[c]}^j + \text{logit}(\hat{\theta}_c^j)\right). \label{eq:updatedprobs}
\end{align}
When the updated estimates are poststratified to the geographic level, they exactly match the ground-truth data. 
When poststratified to other subgroups, the results will generally differ from the uncalibrated MRP results.

For example, researchers studying racially polarized voting may estimate an MRP model, apply the logit shift at the district level to match actual election results, and aggregate adjusted estimates to the racial-group level \citep{Kuriwaki2023}. The resulting race-level vote share estimates will differ from their uncalibrated counterparts.

The logit shift is a minimal, geographic-based adjustment that does not distinguish between individuals within the same geography and has been shown to approximate the posterior distribution of cell-level probabilities conditional on the ground truth \citep{RMO:2023}.



Calibration improves estimates only under certain assumptions \citep[section 3]{RMO:2023}. The shift is uniform within each geography, so if the true error is not uniform, calibration may increase error for some cells even as it decreases average error.
The uniform-shift assumption is more appropriate for homogeneous geographies. Given residential segregation along political and racial lines \citep{Rodden2019}, the logit shift generally works better at smaller scales.